In this work, we conducted a controlled comparison of five LLM agent strategies for navigation in MiniGrid: a reactive baseline, ReAct, Chain of Thought, Tree of Thought, and Buffer of Thought. Across three environments of increasing difficulty, our results show that explicit reasoning or planning can improve navigation performance over purely reactive action selection, especially as the task becomes more challenging. However, these gains are not free as stronger reasoning methods generally incur higher token cost, longer inference time, or initialization overhead. 

We observe that all methods perform strongly on the simplest environment Minigrid-Empty-8x8-v0, suggesting that sophisticated planning is not necessary for easier goal-reaching tasks. As the environment becomes more hazardous and exploration-heavy, the differences between prompting strategies become more pronounced. Chain-of-Thought (CoT) \cite{wei2023chainofthoughtpromptingelicitsreasoning} provides a good middle ground, improving capability while avoiding the overheads associated with Buffer-of-Thought (BoT) and Tree-of-Thought (ToT). Tree-of-Thought \cite{yao2023tree} and Buffer-of-Thought \cite{yang2024buffer} achieve the best performance overall on the harder settings, but they do so with different tradeoffs. ToT is the most expensive in token usage because it evaluates multiple candidate reasoning branches at each step, while BoT can remain comparatively token-efficient but often requires more exploratory steps. By contrast, the reactive baseline remains the cheapest strategy overall in terms of tokens, but it is less reliable once the navigation problem requires longer-horizon reasoning.

Overall, our study suggests that there is no universally best LLM agent strategy for all navigation tasks. The choice method depends on the deployment constraint: if token budget and latency are critical, simpler reactive or CoT-based agents may be preferable; if maximizing task success is the priority, more structured planning methods such as ToT or memory-augmented methods such as BoT become more attractive. Future work can extend this comparison to larger and more capable models, additional MiniGrid \cite{chevalier2023minigrid} tasks involving keys, doors, such as DoorKey, and object manipulation, such as Blocked Unblocked Pickup, and possible hybrid agent designs that may combine deliberative planning with reusable memory. It would also be useful to study methods such as quantization and prompt compression to reduce latency and cost while preserving the benefits of more advanced reasoning strategies.
